\chapter{ヘッジファンドは儲かっているのか}
\label{ch:hfperf}

第 \ref{ch:equityfactor}〜\ref{ch:hft} 章は「戦略」の研究だった。この章は、実際のヘッジファンドの「成績」の研究をまとめる。結論を先に書くと、平均的なヘッジファンドの手数料の後のアルファは小さく、データの偏りを補正するとさらに小さくなる。少数のファンドには本物の腕があるらしいが、それを事前に見分けるのは難しい。

\section{ヘッジファンドのリターンを説明するファクター}

ヘッジファンドは、空売り・レバレッジ・デリバティブを使い、ポジションを頻繁に変えるので、株式の投資信託の評価に使うファクター（市場、規模、割安、モメンタム）では、そのリスクを捉えきれない。

\begin{itemize}
\item \textbf{Fung・Hsieh（1997）}\cite{funghsieh1997}：ヘッジファンドのリターンは、伝統的な資産クラスとの相関が低く、少数の「スタイル」（売買の仕方）の要因で説明できることを示した。
\item \textbf{Fung・Hsieh（2001）}\cite{funghsieh2001}：トレンドフォローの運用者のリターンが、ルックバック・ストラドル（オプション）のリターンに似ることを示した（第 \ref{ch:trend} 章）。
\item \textbf{Agarwal・Naik（2004）}\cite{agarwalnaik2004}：多くの株式のヘッジファンドの戦略のリターンが、株価指数のアウト・オブ・ザ・マネーのプット・オプションを売った形に似る（左の裾のリスクを持つ）ことを示し、オプションの形のファクターを提案した。平均と分散だけの評価は、この裾のリスクを過小評価する。
\item \textbf{Fung・Hsieh（2004）}\cite{funghsieh2004}：株式市場、規模の差、国債の利回りの変化、信用のスプレッドの変化、債券・通貨・商品のトレンドフォロー（ルックバック・ストラドル）の 7 つのファクターのモデルを提案した。このモデルは、ヘッジファンドの指数のリターンの変動の大部分（約 8 割）を説明する。ヘッジファンドの研究の標準のベンチマークになった。
\end{itemize}

\begin{projbox}
「ヘッジファンドの収益の多くは、ルールで再現できるファクター（オルタナティブ・ベータ）で説明できる」という結果（第 \ref{ch:hf} 章の 3 分類の 2）は、この一連の研究に基づく。
\end{projbox}

\section{手数料の後のアルファ}

\subsection{メタ分析}

\begin{paperbox}[Yang・Havranek・Irsova・Novak（2022）\cite{yhin2022}]
2001〜2021 年に発表された 74 の研究から、ヘッジファンドのリターンをリスク・ファクターに回帰した切片（アルファ）の推定値 1,019 個を集めたメタ分析。
\begin{itemize}
\item 報告されたアルファの単純な平均は月 0.36\%（年 4.32\%）。発表の偏り（有意な結果だけが発表されやすい偏り）は小さく、それを補正した推定の多くは月 0.30〜0.40\% の範囲にある。
\item データの偏り（遡って成績が入る偏り、生存者バイアス）を明示的に補正した研究ほど、低いアルファを報告している。
\end{itemize}
\end{paperbox}

Deep Research の答え 3 は、この論文が「データの期間の中央が 2015 年以降の研究では、平均のアルファが負になる」ことを示したと書いたが、照合した 2022 年 7 月の版の要旨と本文の検索では、その記述を見つけられなかった。この報告書では使わない。

\subsection{時期による変化}

\begin{paperbox}[Fung・Hsieh・Naik・Ramadorai（2008）\cite{fhnr2008}]
1995〜2004 年のファンド・オブ・ファンズを分析し、平均的なファンド・オブ・ファンズが有意なアルファを上げたのは、1998 年 10 月〜2000 年 3 月の短い期間だけだったことを示した。アルファを上げたファンドには資金が流入し、流入の後にアルファは下がった。資金の流入が、アルファの機会を食いつぶしていく（資本の形成）。
\end{paperbox}

\subsection{ラグのある市場の感応度}

Asness・Krail・Liew（2001）\cite{akl2001} は、ヘッジファンドの月次のリターンを、同じ月の市場のリターンだけでなく、前の月の市場のリターンにも回帰すると、市場への感応度（ベータ）がずっと大きくなり、アルファの多くが消えることを示した。流動性の低い資産の評価が遅れて市場の動きを反映するため（\ref{sec:smoothing} 節の平滑化）である。Hasanhodzic・Lo（2007）\cite{hasanhodziclo2007} は、いくつかの戦略の種類で、ヘッジファンドのリターンのかなりの部分が、株・債券・通貨・商品・信用などの共通のファクターへの線形の感応度で再現できる（「クローン」が作れる）ことを示した。

\subsection{投資家が実際に得た収益}

\begin{paperbox}[Dichev・Yu（2011）\cite{dichevyu2011}]
ファンドの成績（買って持ち続けた場合の収益）ではなく、投資家の資金の出し入れを考慮した収益（金額で加重した収益）を計算した。投資家は成績の良い後に資金を入れ、悪い後に引き出すので、金額で加重した年率の収益は、買って持ち続けた収益より数\% 低く（SSRN に出た版の要旨では約 4\%）、幅広い株価指数の収益も下回り、変動も大きかった。ヘッジファンドの投資家の実際の経験は、それまで考えられていたよりずっと悪い。

\end{paperbox}

\section{データベースの偏り}

ヘッジファンドには成績の公開の義務がなく、データベースへの報告は任意である。このため、データベースの成績は、実際の平均より良く見える。

\begin{table}[H]
\centering\small
\caption{ヘッジファンドのデータベースの偏り}
\label{tab:hfbias}
\begin{tabular}{@{}P{0.2\linewidth}P{0.42\linewidth}P{0.3\linewidth}@{}}
\toprule
偏り & 仕組み & 主な研究 \\
\midrule
生存者バイアス & 閉じた・報告をやめたファンドがデータから消える。 & Liang（2000）\cite{liang2000}、Brown・Goetzmann・Ibbotson（1999）\cite{bgi1999}、Fung・Hsieh（2000）\cite{funghsieh2000} \\
遡って入る偏り & 新しいファンドは、成績が良かったときだけデータベースに加わり、過去の成績も一緒に入る（インスタント・ヒストリー）。 & Fung・Hsieh（2000）、Malkiel・Saha（2005）\cite{malkielsaha2005} \\
自己選択 & 報告を始めるのは成績が良い後、やめるのは悪い後。 & Agarwal・Fos・Jiang（2013）\cite{afj2013}、Aiken・Clifford・Ellis（2013）\cite{ace2013} \\
評価の平滑化・操作 & 流動性の低い資産の評価をならす。小さな損失を報告したがらない。 & Getmansky・Lo・Makarov（2004）\cite{glm2004}、Bollen・Pool（2009）\cite{bollenpool2009} \\
\bottomrule
\end{tabular}
\end{table}

偏りの大きさの推定は、データベース・期間・方法で違う。Yang ほか（2022）のまとめによれば、生存者バイアスの推定は年 2\% 前後から 3\% 台（例えば、ある研究は年 1.8〜2.4\%、別の研究は年 2.0〜3.6\%）で、小さく若くレバレッジの高いファンドではさらに大きい。Liang（2000）\cite{liang2000} は年 2\% を超える生存者バイアスを報告した。Malkiel・Saha（2005）\cite{malkielsaha2005} は、遡って入る偏りと生存者バイアスがともに大きく、ヘッジファンドのリターンは一般に思われているより低く、リスクは高いと結論した。

\begin{itemize}
\item \textbf{報告をやめたファンド}：Aiken・Clifford・Ellis（2013）\cite{ace2013} は、ファンド・オブ・ファンズの保有の開示から、データベースに報告していないファンドの成績を調べ、データベースから抜けたファンドのその後の成績が非常に悪いことを示した（成績の分布の左の裾がデータから消えている）。
\item \textbf{ゼロの前後の不自然さ}：Bollen・Pool（2009）\cite{bollenpool2009} は、ヘッジファンドの月次のリターンの分布で、小さな損失が小さな利益より不自然に少ないこと（ゼロでの不連続）を示し、損失の報告を避ける評価の操作の可能性を指摘した。
\end{itemize}

\begin{projbox}
LLM-BLM の小型株のデータの作り方（生き残った会社だけを選び、その偏りを SEC のデータで別に見積もった）は、ヘッジファンドのデータベースの生存者バイアスと同じ問題への対処である。偏りは比べる相手との比較で一部打ち消せるが、選び方そのものの偏りは残る。
\end{projbox}

\section{腕は続くか}

\begin{itemize}
\item \textbf{Kosowski・Naik・Teo（2007）}\cite{knt2007}：ブートストラップとベイズの方法で、上位のヘッジファンドのアルファが運だけでは説明できないこと、1 年の単位で成績の持続があることを示した。
\item \textbf{Jagannathan・Malakhov・Novikov（2010）}\cite{jmn2010}：データベースの偏りと推定の誤差を考慮しても、優れたファンドの成績は 3 年の単位で持続し（「熱い手」）、劣ったファンドの持続は弱いことを示した。
\item \textbf{Sun・Wang・Zheng}\cite{swz2018}：ヘッジファンドの業界全体が不調だった時期の成績（DownsideReturns）で並べると、上位の 5 分の 1 は下位の 5 分の 1 を、リスクを調整した翌年の成績で約 7\% 上回り、その予測力は 3 年まで続いた。好調だった時期の成績で並べても、その後の差は出なかった。上げ相場では、リスクを取っただけのファンドも良く見えるので、腕と運を見分けにくい（Deep Research の答え 3 は、この論文の題名を誤って書いていた。付録~\ref{ch:verify}）。
\item \textbf{Titman・Tiu（2011）}\cite{titmantiu2011}：ファクターとの連動の小さい（決定係数の低い）ファンドほど、シャープレシオとアルファが高く、手数料も高いことを示した。腕のある運用者は、ファクターのリスクを打ち消し、固有の収益を売る。
\end{itemize}

\section{規模と成績}

\begin{itemize}
\item \textbf{Berk・Green（2004）}\cite{berkgreen2004}：投資信託の理論で、腕のある運用者には資金が集まり、規模が大きくなるとアルファが下がる（規模の不経済）ので、均衡では投資家の手数料の後のアルファはゼロになり、成績の持続も見えなくなることを示した。腕は存在しても、その価値は運用者が手数料として取る。
\item \textbf{Pástor・Stambaugh・Taylor（2015）}\cite{pst2015}：投資信託で、運用会社の規模より、業界全体の規模の拡大がアルファを下げること（業界の規模の不経済）を示した。運用者の腕は時代とともに上がっているが、業界の拡大がそれを打ち消している。
\item \textbf{Getmansky（2012）}\cite{getmansky2012}：ヘッジファンドの成績と規模の関係は凹で、最適な規模があることを示した。
\item \textbf{Yin（2016）}\cite{yin2016}：ヘッジファンドの成績は規模とともに下がるが、運用者の報酬（管理報酬は規模に比例）は規模とともに増えるので、運用者は投資家にとって最適な規模を超えて資金を集める誘因を持つことを示した。
\end{itemize}

\section{手数料は誰のものか}

\begin{paperbox}[Ben-David・Birru・Rossi（2020）「成功報酬の成績」\cite{bbr2020}]
1995〜2016 年の約 22 年間のヘッジファンドで、実際に支払われた成功報酬は、利益の合計に対して、契約上の平均（約 20\%）の約 2.5 倍、つまり約 50\% に達した。投資家は、安全な金利を超えて得た収益（リスクの調整の前）の 1 ドルのうち、約 36 セントしか手にしなかった。原因は、(1) 成功報酬は利益が出た年に払われ、後の損失で返されない（契約の非対称）、(2) 投資家が成績を追って資金を出し入れする、(3) 損をしてハイウォーターマークの下に沈んだファンドが閉じる（その後の回復で報酬が相殺されない）、ことである。ファンドの生涯の成績と、受け取った成功報酬の関係は弱い。
\end{paperbox}

\begin{figure}[H]
\centering
\begin{tikzpicture}[x=1.2mm,y=1mm]
\fill[c1] (0,0) rectangle (36,10); \node[font=\footnotesize,white] at (18,5) {投資家 36 セント};
\fill[c2!80] (36,0) rectangle (100,10); \node[font=\footnotesize,white] at (68,5) {運用者（管理報酬＋成功報酬）64 セント};
\draw[ink2] (0,0) rectangle (100,10);
\node[lab,anchor=north west] at (0,-1) {安全な金利を超えた、手数料の前の収益 1 ドル（1995〜2016 年、リスクの調整の前）};
\end{tikzpicture}
\caption{ヘッジファンドの収益の分け前（Ben-David・Birru・Rossi の NBER の要旨の値から描いた）。}
\label{fig:feesplit}
\end{figure}

\ref{sec:fees} 節の手数料の例（手数料の前の収益の 36\% が運用者に渡る）は 1 年の計算だった。複数年・多数のファンド・投資家の行動を考えると、運用者の取り分はさらに大きくなる。

\section{まとめ}

\begin{keybox}[ヘッジファンドの成績について確立していること]
\begin{itemize}
\item ヘッジファンドのリターンの大部分は、オプションの形を含むファクター（Fung・Hsieh の 7 ファクターなど）で説明できる。
\item 報告されたアルファの平均は月 0.3〜0.4\% 程度だが、データの偏りを補正すると小さくなる。データベースの偏り（生存者・遡り・自己選択）は年に数\% の大きさになりうる。
\item 投資家が実際に得た収益は、資金の出し入れのタイミングのため、ファンドの成績よりさらに低い。
\item 手数料は、手数料の前の超過収益の大部分を取る。
\item 少数の優れたファンドの成績には持続があるが、規模の拡大と資金の流入がアルファを小さくする。
\end{itemize}
\end{keybox}

LLM-BLM の研究が時価総額加重に勝てなかったことは、プロの世界の平均的な成績から見て、驚くべきことではない。プロの多くも、手数料とコストの後では、市場と同じリスクのファクターの組み合わせに勝てていない。
